\chapter{Controle por modos deslizantes adaptativo (ASMC)}
\label{cap:smc}

\section{Lei de controle}
Arquivo \code{control/sliding\_mode\_controller.cpp}:
\begin{align}
  \vect s &= \ev + \lambda\,\ep && \text{(superfície de deslizamento)}\\
  \vect u &= -\lambda\,\ev - k_s\,\vect s - \hat k\,\sat(\vect s/\phi)
            && \text{(lei de alcance com camada-limite)}\\
  \dot{\hat k} &= \gamma\left(\lVert\vect s\rVert - \sigma\hat k\right)\ \text{se } \lVert\vect s\rVert>\phi,
  \quad \hat k\in[k_{min},k_{max}] && \text{(adaptação com $\sigma$-modificação)}
\end{align}
em que $\sat(\vect s/\phi) = \vect s/\max(\lVert\vect s\rVert,\phi)$ (saturação vetorial).
\code{yaml}: $\lambda=0{,}8$, $k_s=1$, $\hat k(0)=0{,}5$, $k_{max}=3$, $\gamma=0{,}5$,
$\sigma=0{,}1$, $\phi=\SI{0,3}{m/s}$.

\section{Projeto da superfície}
Sobre $\vect s=0$ vale $\ev=-\lambda\ep$ e, como $\dot{\vect e}_p=\ev$,
\begin{equation}
  \dot{\vect e}_p = -\lambda\,\ep \quad\Longrightarrow\quad \ep(t)=\ep(0)e^{-\lambda t}.
\end{equation}
Logo $\lambda$ é a \textbf{taxa de convergência na fase de deslizamento}, com
$\tau=1/\lambda=\SI{1,25}{s}$ e acomodação $\approx 4/\lambda=\SI{5}{s}$ no \code{yaml}.

\paragraph{Superfície ótima.} Tratando $e_v$ como a ``entrada'' da dinâmica
$\dot e_p=e_v$ e minimizando $\int(q_p e_p^2 + q_v e_v^2)dt$, o LQR escalar dá
$e_v=-\sqrt{q_p/q_v}\,e_p$, isto é,
\begin{equation}
  \boxed{\lambda^\star = \sqrt{q_p/q_v}}
\end{equation}
(Utkin). Com os pesos do LQR ($q_p=10$, $q_v=2$): $\lambda^\star=2{,}24$. Esse é um
critério objetivo para a comparação: mesma ponderação que o LQR/MPC.

\section{Fase de alcance e estabilidade (Lyapunov)}
Derivando $\vect s$ com $\dot{\vect e}_v=\vect u-\vect a_t$:
\begin{equation}
  \dot{\vect s} = \vect u - \vect a_t + \lambda\ev
  = -k_s\vect s - \hat k\,\sat(\vect s/\phi) - \vect a_t .
\end{equation}
Com $V=\tfrac12\lVert\vect s\rVert^2$, fora da camada-limite:
\begin{equation}
  \dot V = -k_s\lVert\vect s\rVert^2 - \hat k\lVert\vect s\rVert - \vect s\T\vect a_t
  \le -k_s\lVert\vect s\rVert^2 - (\hat k - a_{t,max})\lVert\vect s\rVert .
\end{equation}
Basta $\hat k > a_{t,max}$ para alcançar a superfície em tempo finito,
$t_r\le\lVert\vect s(0)\rVert/(\hat k-a_{t,max})$. A adaptação faz $\hat k$ crescer até
dominar a aceleração do alvo sem conhecer $a_{t,max}$ a priori. A
$\sigma$-modificação impede a deriva de $\hat k$ com ruído.

\section{Camada-limite = PD de alto ganho}
Dentro de $\lVert\vect s\rVert<\phi$, $\sat(\vect s/\phi)=\vect s/\phi$ e a lei fica
\emph{linear}:
\begin{equation}
  u = -\lambda e_v - k_{eq}(e_v+\lambda e_p),\qquad k_{eq}=k_s+\frac{\hat k}{\phi},
  \qquad\Longrightarrow\qquad
  k_{p,eq} = \lambda k_{eq},\quad k_{d,eq}=\lambda + k_{eq}.
  \label{eq:smc-pd}
\end{equation}
Com o \code{yaml}:
\begin{center}
\begin{tabular}{@{}lccc@{}}
\toprule
$\hat k$ & $k_{eq}$ & $k_{p,eq}$ & $k_{d,eq}$ \\
\midrule
0,5 (inicial) & 2,67 & 2,13 & 3,47 \\
1,5 & 6,0 & 4,8 & 6,8 \\
3,0 ($k_{max}$) & 11,0 & 8,8 & 11,8 \\
\bottomrule
\end{tabular}
\end{center}
\begin{atencao}
\eqref{eq:smc-pd} permite usar \emph{todas} as ferramentas lineares no SMC. Um ponto
contraintuitivo: como $e_v$ é a velocidade de \emph{referência} (\cref{eq:ceq}), o
controlador equivalente é $C_{eq}(s)=k_{p,eq}s/(s+k_{d,eq})$. Quando $k_{eq}\to\infty$,
$C_{eq}\to\lambda s$ para $\omega\ll k_{eq}$: a malha tende a $L\approx\lambda T_v(s)/s$
(a referência ``desliza'' com $v_{ref}\approx-\lambda e_p$). Por isso, com PX4 e
\SI{0,1}{s} de atraso, a margem de fase \emph{não} cai com $\hat k$: ela fica em
\ang{72}, \ang{74} e \ang{75} para $\hat k=0{,}5$, $1{,}5$ e $3$
(\code{projeto\_smc.m}). O custo de um $\hat k/\phi$ alto está em outro lugar:
\begin{itemize}
  \item o ganho sobre o \emph{ruído} de $e_p$ cresce como $k_{p,eq}=\lambda k_{eq}$
        (8,8 com $k_{max}$): a aceleração comandada fica ruidosa;
  \item fora da camada-limite o comando satura em \code{max\_acc} e a malha deixa de
        ser linear.
\end{itemize}
Escolha $k_{max}/\phi$ pelo nível de ruído aceitável em $\vect a$, não pela margem.
\end{atencao}

\section{Tempo discreto e chattering}
Em tempo discreto a trajetória não fica sobre $\vect s=0$: ela oscila numa faixa
(\emph{modo quase-deslizante}) de largura $\approx\hat k\,\Ts$. Com atraso $\tau$, a
faixa vira $\approx \hat k(\Ts+\tau)$. Para não haver chaveamento (chattering), escolha
\begin{equation}
  \phi \gtrsim \hat k_{max}\,(\Ts + \tau) = 3\cdot(0{,}05+0{,}1) = \SI{0,45}{m/s}.
\end{equation}
O \code{yaml} usa $\phi=0{,}3$. Com \SI{0,1}{s} de atraso, vale testar $\phi=0{,}45$ ou
reduzir $k_{max}$. Uma alternativa sem camada-limite é a lei de alcance discreta de
Gao, $s_{k+1}=(1-q\Ts)s_k-\varepsilon\Ts\sgn(s_k)$, com $1-q\Ts>0$.

\section{Ajuste ótimo}
\begin{enumerate}
  \item $\lambda$ pela superfície ótima ($\lambda^\star=\sqrt{q_p/q_v}$) ou por $t_s$.
  \item $k_s$ e $\phi$ pelo PD equivalente: confira a PM no pior caso e limite
        $k_{p,eq}=\lambda(k_s+k_{max}/\phi)$ pelo ruído aceitável na aceleração.
  \item $k_{max}$ $\approx$ 1,5$\times$ a aceleração máxima esperada do alvo.
  \item $\gamma,\sigma$ por simulação: $\gamma$ maior adapta mais rápido e fica mais
        sensível a ruído.
  \item Refinamento: \code{otimizar\_ganhos.m} com $\theta=[\lambda,k_s,\phi]$.
\end{enumerate}
